\documentclass[11pt]{article}
\usepackage[margin=1in]{geometry}
\usepackage{amsmath,amssymb}
\usepackage{booktabs}
\usepackage{hyperref}

\title{Do contrastive embeddings converge to a log-probability?\\
\large InfoNCE vs.\ LeCun's MSE\,+\,SIGReg (LeJEPA)}
\author{}
\date{}

\begin{document}
\maketitle

\section{Question}
Given a contrastive-style encoder trained to its optimum, does the distance
between two (class) centers in embedding space converge to a log-probability?
We compare two regimes: (i) InfoNCE / SimCLR-style losses, and (ii) LeCun's
MSE\,+\,SIGReg objective (LeJEPA,
\href{https://arxiv.org/abs/2511.08544}{arXiv:2511.08544}).

\section{InfoNCE: the log-density-ratio comes from the partition function}
For anchor $x$, positive $x^+$, temperature $\tau$, and critic $f$,
\begin{equation}
\mathcal{L}_{\mathrm{InfoNCE}}
= -\frac{f(x,x^+)}{\tau}
+ \log\sum_{k}\exp\!\frac{f(x,x_k)}{\tau}.
\end{equation}
With unbounded negatives and sufficient capacity, the minimizing critic is the
pointwise mutual information (Oord et al.; Wang \& Isola):
\begin{equation}
\frac{f^\star(x,x^+)}{\tau}
= \log\frac{p(x^+\mid x)}{p(x^+)} + \text{const}
= \operatorname{PMI}(x,x^+).
\label{eq:pmi}
\end{equation}
The \emph{softmax denominator} is what forces the critic onto a log-ratio:
InfoNCE is implicitly a density-ratio estimator. On the unit sphere the
\emph{similarity} is what lands on a log-probability, and distance is affine in
it:
\begin{equation}
d^2(z_i,z_k) = 2 - 2\cos\theta_{ik}
= 2 - 2\tau\,\operatorname{PMI}(x_i,x_k)/\text{scale}.
\end{equation}

\paragraph{Supervised case.}
With supervised positives (same-class $x^+$), $p(x^+\mid x)$ is the same-class
conditional, and the center-to-center object becomes a
\emph{log-likelihood-ratio} between class conditionals $p(\cdot\mid a)$ and
$p(\cdot\mid b)$ --- a measure of class confusability/overlap.

\paragraph{Two caveats.}
\begin{enumerate}
\item \textbf{Non-collapse required.} If the embedding dimension is too small
(and $\tau$ small), the supervised optimum is a simplex equiangular tight frame
(neural collapse): all class centers equidistant, and the likelihood
information is destroyed.
\item \textbf{Affine, and a ratio.} The readout is affine in a log-density
\emph{ratio}, with a $\tau$-dependent scale and an unidentified additive
constant --- not a calibrated log-likelihood.
\end{enumerate}

\section{LeJEPA (MSE\,+\,SIGReg): the mechanism is removed by design}
The objective has no partition function, no negatives, no critic:
\begin{equation}
\mathcal{L}
= \underbrace{\bigl\|z_{\mathrm{pred}} - z_{\mathrm{target}}\bigr\|^2}_{\text{alignment}}
+ \lambda\,
\underbrace{\mathcal{L}_{\mathrm{SIGReg}}}_{\text{marginal}\,\to\,\mathcal{N}(0,I)},
\qquad
\mathcal{L}_{\mathrm{SIGReg}}
= \sum_k \bigl|\phi_{\mathrm{emp}}(\omega_k) - \phi_{\mathcal{N}}(\omega_k)\bigr|^2,
\end{equation}
where SIGReg is an Epps--Pulley characteristic-function test pinning the
\emph{aggregate} embedding marginal to an isotropic Gaussian. Because there is
no softmax denominator, \emph{the InfoNCE$\to$PMI argument does not apply}:
nothing drives the pairwise distance to a log-likelihood-ratio.

\paragraph{What is a log-probability here.}
The constraint makes the \emph{norm}, not the pairwise distance, affine in a
log-density --- and it is the density under the \emph{imposed prior}, not the
data:
\begin{equation}
\log p_{\mathrm{ref}}(z) = -\tfrac12\|z\|^2 + \text{const}
\quad\Longrightarrow\quad
\|z\|^2 = -2\log p_{\mathrm{ref}}(z).
\end{equation}
Recovering the true \emph{data} log-likelihood would additionally require the
encoder's change-of-variables Jacobian (as in a normalizing flow); SIGReg
constrains only the marginal, not the Jacobian, so data log-likelihood is
\emph{not} recovered.

\paragraph{Supervised / class centers.}
Supervised MSE (pull same-class together) $+$ SIGReg (spread the set to
$\mathcal{N}(0,I)$) yields Gaussian class blobs packed so the mixture
$\approx\mathcal{N}(0,I)$. Then the class-center distance is a
Mahalanobis/Gaussian separation, which equals a log-likelihood-ratio
\emph{only under an added LDA-style assumption}: each class conditional Gaussian
with shared covariance. SIGReg guarantees only that the \emph{aggregate} is
Gaussian, not the per-class conditionals. With balanced classes the centers
pack near-symmetrically, so distances reflect class count/priors (packing
geometry), not confusability --- the Gaussian-packing analogue of the ETF
collapse.

\section{Summary}
\begin{center}
\begin{tabular}{@{}lll@{}}
\toprule
 & InfoNCE / SimCLR & LeJEPA (MSE\,+\,SIGReg)\\
\midrule
Similarity $\to$ & $\log\frac{p(x^+\mid x)}{p(x^+)}$ (PMI) & nothing (no critic)\\
Pinned to a log-prob & pairwise distance (affine in $-$log-ratio) & the \emph{norm}: $\|z\|^2=-2\log p_{\mathrm{ref}}$\\
Center dist.\ $=$ log-lik.\ ratio? & yes, up to affine $+$ non-collapse & only under equal-cov.\ Gaussian assumption\\
\bottomrule
\end{tabular}
\end{center}

\noindent
\textbf{Bottom line.} InfoNCE's distance encodes an \emph{estimated}
log-density-ratio via its partition function. LeCun's MSE\,+\,SIGReg deliberately
replaces that with an \emph{imposed} marginal prior: you gain provable
distributional control and lose the automatic ``distance $=$ log-likelihood''
readout. To make class-center distances into log-likelihood-ratios under SIGReg,
you must add a per-class second-moment constraint; the aggregate Gaussian alone
will not do it.

\section{An NPLM-based contrastive loss}
The quantity we want --- a \emph{calibrated} log-likelihood-ratio --- is exactly
what the New Physics Learning Machine (NPLM, D'Agnolo \& Wulzer) estimates by
construction, so it is a natural loss to build a contrastive objective on.

\paragraph{Reframing.} Both SimCLR and the supervised question are asking for the
pointwise mutual information
\begin{equation}
\operatorname{PMI}(x,x') = \log\frac{p(x,x')}{p(x)\,p(x')},
\end{equation}
where the joint $p(x,x')$ is the distribution of \emph{positive pairs}
(augmentations of the same point, or same-class pairs) and the reference
$p(x)p(x')$ is \emph{independent pairs} (a shuffled batch). InfoNCE estimates
this ratio with a coupled softmax; NPLM estimates the same ratio with a
maximum-likelihood (Neyman--Pearson) loss.

\paragraph{The loss.} NPLM parametrizes the alternative density as
$n(x)=n_{\mathrm{ref}}(x)\,e^{f(x)}$ and minimizes
$\mathcal{L}[f]=\sum_{\mathrm{ref}}(e^{f}-1)-\sum_{\mathrm{data}}f$. Ported to
pairs with a critic $g(x,x')$,
\begin{equation}
\mathcal{L}[g]
= \mathbb{E}_{(x,x')\sim p(x)p(x')}\!\big[e^{g(x,x')}-1\big]
- \mathbb{E}_{(x,x')\sim p(x,x')}\big[g(x,x')\big].
\label{eq:nplm}
\end{equation}
The stationarity condition $p(x)p(x')\,e^{g^\star}=p(x,x')$ gives
\begin{equation}
g^\star(x,x') = \log\frac{p(x,x')}{p(x)p(x')} = \operatorname{PMI}(x,x').
\end{equation}
Crucially the $e^{g}-1$ term forces $\mathbb{E}_{\mathrm{ref}}[e^{g^\star}]=1$,
which pins the additive constant: unlike InfoNCE, the recovered value is the
\emph{absolute}, calibrated log-likelihood-ratio.

\paragraph{Distance geometry.} Choosing
\begin{equation}
g(x,x') = -\tfrac12\|z(x)-z(x')\|^2 + b(x) + b(x'),
\end{equation}
with a learned log-density bias $b$, the optimum yields
\begin{equation}
\|z(x)-z(x')\|^2 = -2\operatorname{PMI}(x,x') + 2b(x) + 2b(x'),
\end{equation}
so squared distance is exactly affine in the negative \emph{calibrated}
log-likelihood-ratio, with $b$ absorbing the per-point marginal density.

\paragraph{Test statistic.} $t=-2\,\mathcal{L}_{\min}$ is an asymptotic
log-likelihood-ratio test (Wilks) --- the same object trained is the
signal-detection statistic, which is the point of NPLM.

\paragraph{Hybrid with LeJEPA.} Keep SIGReg to pin the marginal
$p(z)=\mathcal{N}(0,I)$ (so $\|z\|^2=-2\log p_{\mathrm{ref}}$) \emph{and} add the
NPLM ratio head~\eqref{eq:nplm} for calibrated pairwise log-likelihood-ratios.
This gives both a controlled marginal geometry and pairwise distances that are
genuine calibrated log-likelihood-ratios --- the ``distance $=$ log-probability''
property neither loss delivers alone.

\subsection{Pros and cons}
\paragraph{Pros.}
\begin{itemize}
\item \textbf{Calibration.} The $e^{g}-1$ term fixes the additive constant, so
the critic recovers the \emph{absolute} log-likelihood-ratio, not
``$\operatorname{PMI}+\text{const}$'' as in InfoNCE. This is the ``distance $=$
log-likelihood'' readout, genuinely calibrated.
\item \textbf{No coupled partition function.} The reference expectation is
separable; there is no softmax denominator locking many negatives together. One
needs enough shuffled pairs, not a large coupled negative set.
\item \textbf{Built-in test statistic.} $t=-2\mathcal{L}_{\min}$ is an asymptotic
log-likelihood-ratio test, directly usable for signal detection --- the trained
object and the test statistic coincide.
\item \textbf{Weights and imbalance.} Class weights / imbalance enter as simple
sample weights, without reweighting a normalization constant.
\item \textbf{Interpretable geometry.} With the distance parametrization, pairwise
distances and class-center separations are calibrated log-likelihood-ratios of
confusability, free of the ETF-collapse scale ambiguity.
\end{itemize}

\paragraph{Cons.}
\begin{itemize}
\item \textbf{Exponential instability.} $e^{g}$ diverges when $g$ is
large-positive on reference pairs --- the classic NPLM training pathology.
Requires weight clipping / bounding $\|g\|$, a kernel version (Falkon-NPLM), or
regularization to control the variance of $\mathbb{E}_{\mathrm{ref}}[e^{g}]$.
\item \textbf{Reference-sample cost.} It remains density-ratio estimation; the
calibration is only as good as the shuffled-pair coverage of $p(x)p(x')$.
\item \textbf{Identifiability of the distance form.} The bias $b$ must be
expressive enough to represent marginal-density variation, or the distance term
is forced to explain structure it cannot represent.
\item \textbf{No automatic marginal control.} Unlike SIGReg, the loss says
nothing about the shape of $p(z)$; if a controlled marginal is wanted it must be
added (hence the hybrid above).
\item \textbf{Higher-variance gradients.} The exponential term generally yields
noisier gradients than the log-softmax of InfoNCE, so it can need smaller steps /
more careful optimization.
\end{itemize}

\section{Building an embedded log-likelihood space (NPLM $+$ SIGReg)}
The aim of this section is a concrete objective, structured like InfoNCE (an
encoder $z=h_\theta(x)$, positive pairs, in-batch negatives), whose optimum makes
the \emph{embedding geometry itself} the log-likelihood surface --- not up to an
unidentified constant, but ``in the true sense.''

\subsection{Design target: what ``true log-likelihood space'' means}
We ask the encoder $z=h_\theta(x)\in\mathbb{R}^d$ to satisfy two readouts
simultaneously, against a fixed reference $q(z)=\mathcal{N}(0,I)$:
\begin{align}
\text{(marginal)}\quad & \log p(x) = -\tfrac12\|z\|^2 + c_0,
\label{eq:marg}\\[2pt]
\text{(interaction)}\quad & \log\frac{p(x,x')}{p(x)p(x')} = \langle z, z'\rangle.
\label{eq:inter}
\end{align}
Adding \eqref{eq:marg} for $x$ and $x'$ to \eqref{eq:inter} gives the \emph{joint}
log-density as a single quadratic,
\begin{equation}
\log p(x,x')
= -\tfrac12\|z\|^2 - \tfrac12\|z'\|^2 + \langle z,z'\rangle + c_1
= -\tfrac12\|z-z'\|^2 + c_1.
\label{eq:joint}
\end{equation}
So the design target is precisely: \emph{marginals standard Gaussian, positive
pairs Gaussian in the embedding difference}. Every likelihood object --- marginal,
conditional, ratio, joint --- is then a quadratic in $z$, which is what we mean by
the embedding \emph{being} log-likelihood space. Equations \eqref{eq:marg} and
\eqref{eq:inter} are supplied by two different terms: SIGReg delivers the
marginal, NPLM delivers the interaction.

\subsection{The embedded NPLM critic and its minimizer}
Take the InfoNCE-style critic --- inner product plus a learned per-point bias ---
\begin{equation}
g_\theta(x,x') = \langle z, z'\rangle + b_\theta(x) + b_\theta(x'),
\qquad z=h_\theta(x),
\label{eq:critic}
\end{equation}
and plug it into the NPLM density-ratio loss between the joint (positives) and the
product of marginals (independent pairs),
\begin{equation}
\mathcal{L}_{\mathrm{NPLM}}[g]
= \mathbb{E}_{p(x)p(x')}\!\big[e^{g}-1\big]
- \mathbb{E}_{p(x,x')}\big[g\big].
\label{eq:Lnplm}
\end{equation}
Its unconstrained functional derivative,
$\delta\mathcal{L}/\delta g = p(x)p(x')e^{g}-p(x,x')$, vanishes at
\begin{equation}
g^\star(x,x') = \log\frac{p(x,x')}{p(x)p(x')} = \operatorname{PMI}(x,x').
\end{equation}
For the \emph{constrained} family \eqref{eq:critic} the minimizer is a projection:
using $p(x,x')=p(x)p(x')e^{g^\star}$,
\begin{equation}
\mathcal{L}_{\mathrm{NPLM}}[g]-\mathcal{L}_{\mathrm{NPLM}}[g^\star]
= \mathbb{E}_{p(x)p(x')}\!\Big[e^{g}-e^{g^\star}-e^{g^\star}\big(g-g^\star\big)\Big]
= D_\phi\!\big(g^\star\,\|\,g\big)\ \ge 0,
\label{eq:bregman}
\end{equation}
the Bregman divergence generated by the convex $\phi(t)=e^{t}$ under the reference
measure. Hence the encoder learns the \emph{best exponential-Bregman approximation
of $\operatorname{PMI}$} representable by \eqref{eq:critic}; if $\operatorname{PMI}$
is genuinely bilinear-plus-separable, the fit is exact and
$\langle z,z'\rangle+b(x)+b(x') = \operatorname{PMI}(x,x')$.

\subsection{The residual gauge, and how SIGReg removes it}
Equation \eqref{eq:critic} has exactly the gauge freedom that leaves InfoNCE's
constant unidentified:
\begin{equation}
z \mapsto Rz + a,\qquad b(x)\mapsto b(x) - \langle a, z\rangle - \tfrac12\|a\|^2 + \beta,
\end{equation}
for any rotation $R$, shift $a$, constant $\beta$ --- all leave $g$ (hence
$\mathcal{L}_{\mathrm{NPLM}}$) invariant. So NPLM alone fixes the \emph{interaction}
\eqref{eq:inter} but not the split between $z$ and $b$: the map to
``log-likelihood space'' is undetermined. Impose SIGReg on the marginal,
\begin{equation}
\mathcal{L}_{\mathrm{SIGReg}}
= \sum_k\big|\phi_{\mathrm{emp}}(\omega_k)-\phi_{\mathcal{N}}(\omega_k)\big|^2
\ \xrightarrow{\ \min\ }\ p(z)=\mathcal{N}(0,I).
\end{equation}
This kills the shift/rotation freedom (only $R$ with $a=0$ survives, an irrelevant
frame choice) and, via \eqref{eq:marg}, \emph{identifies the bias as the marginal
log-density},
\begin{equation}
b_\theta(x)\ \longrightarrow\ -\tfrac12\|z\|^2 + \text{const}.
\end{equation}
Substituting back into \eqref{eq:critic} collapses the critic to the pure distance
form $g=-\tfrac12\|z-z'\|^2+\text{const}$, and \eqref{eq:joint} holds. The two
terms are therefore complementary by design: \textbf{NPLM pins the interaction and
its absolute scale ($\mathbb{E}_{\mathrm{ref}}[e^{g^\star}]=1$), SIGReg pins the
marginal and removes the residual gauge.} Neither alone yields
\eqref{eq:marg}--\eqref{eq:joint}; together they do.

\subsection{The combined objective and its minibatch estimator}
\begin{equation}
\boxed{\;
\mathcal{L}
= \underbrace{\mathbb{E}_{p(x)p(x')}\!\big[e^{g_\theta}-1\big]
- \mathbb{E}_{p(x,x')}\big[g_\theta\big]}_{\text{NPLM interaction}}
\;+\;\lambda\,
\underbrace{\sum_k\big|\phi_{\mathrm{emp}}(\omega_k)-\phi_{\mathcal{N}}(\omega_k)\big|^2}_{\text{SIGReg marginal}}
\;}
\end{equation}
For a batch $\{(x_i,x_i^+)\}_{i=1}^N$ with $g_{ij}=g_\theta(x_i,x_j^+)$, use the
diagonal as positives and the off-diagonal (shuffled) pairs as the
product-of-marginals reference:
\begin{equation}
\widehat{\mathcal{L}}_{\mathrm{NPLM}}
= \frac{1}{N(N-1)}\sum_{i\neq j}\big(e^{g_{ij}}-1\big)
- \frac{1}{N}\sum_i g_{ii}.
\end{equation}
This reuses \emph{exactly} the InfoNCE batch geometry; the only change is the loss
functional:
\begin{equation}
\widehat{\mathcal{L}}_{\mathrm{InfoNCE}}
= -\frac1N\sum_i \log\frac{e^{g_{ii}}}{\sum_j e^{g_{ij}}}
\qquad\text{vs.}\qquad
\widehat{\mathcal{L}}_{\mathrm{NPLM}}
= \frac{1}{N(N-1)}\sum_{i\neq j}e^{g_{ij}} - \frac1N\sum_i g_{ii} - 1.
\end{equation}
InfoNCE is the \emph{self-normalized} (per-anchor softmax) estimator of the same
log-ratio; NPLM is the \emph{un-normalized maximum-likelihood} estimator, which is
why it retains the absolute normalization that InfoNCE discards. Reading $g_{ii}$
off the diagonal as $-\tfrac12\|z_i-z_i^+\|^2$ (once SIGReg is active) makes the
positive term a plain squared-distance pull, and the reference term a repulsion
weighted by $e^{-\frac12\|z_i-z_j^+\|^2}$.

\subsection{Result}
At the joint optimum the embedding satisfies the three quadratic identities
\begin{equation}
\log p(x) = -\tfrac12\|z\|^2 + c_0,\quad
\log\frac{p(x,x')}{p(x)p(x')} = \langle z,z'\rangle,\quad
\log p(x,x') = -\tfrac12\|z-z'\|^2 + c_1,
\end{equation}
so marginal densities, pairwise log-likelihood-ratios, and joint log-densities are
all read directly off the geometry --- an embedding that approximates
log-likelihood space in the true (calibrated, gauge-fixed) sense. The
approximation is exact when $\operatorname{PMI}$ lies in the bilinear-plus-separable
family \eqref{eq:critic}; otherwise \eqref{eq:bregman} says the encoder returns the
closest such surface in exponential-Bregman divergence. The practical caveats of
\S5.1 (exponential instability of $e^{g}$, reference-pair coverage) carry over and
set the sole implementation burden.

\section{The exp-34e hybrid, and a configurable NPLM geometry}
The SuperSig codebase already contains a hybrid of the contrastive and SIGReg
families (\texttt{experiments/34e\_hybrid\_contrastive.py}, via
\texttt{train\_simclr\_sigreg}). This section places it in the framework above,
identifies why it stops short of the \S6 identities, and shows that a single loss
class --- with the NPLM interaction as one option --- spans it and many neighbours.

\subsection{What exp-34e does, in this language}
For two augmented views $v_1,v_2$ with raw embeddings $z=h_\theta(v)$, the hybrid
half optimizes
\begin{equation}
\mathcal{L}_{\text{34e}}
= \underbrace{\mathrm{NTXent}\!\big(\hat z,\hat z'\big)}_{\text{instance positives, }\hat z = z/\|z\|}
+ \lambda\,\underbrace{\mathrm{SIGReg}(z)}_{\text{raw }z\,\to\,\mathcal{N}(0,I)},
\qquad \hat z = z/\|z\|.
\end{equation}
In the vocabulary above: NT-Xent is the \emph{self-normalized} (softmax) estimator
of $\operatorname{PMI}$ (\S5, \S6.4) with \emph{instance} positives, and SIGReg
supplies the marginal (\S3). The supervised twin \texttt{train\_supcon\_sigreg}
just swaps instance positives for same-label positives. So exp-34e is already the
``InfoNCE $+$ SIGReg'' hybrid of \S5 --- with one consequential detail.

\subsection{Why it stops short: two geometries}
NT-Xent acts on the \emph{normalized} embedding $\hat z$ (angles on the unit
sphere; scale-invariant), whereas SIGReg acts on the \emph{raw} embedding $z$ (it
is entirely about the radial and coordinate distribution). The interaction term is
blind to $\|z\|$; the marginal term is largely about $\|z\|$. The two therefore
speak different geometries and are only loosely coupled: one obtains features whose
\emph{angles} encode $\operatorname{PMI}+\text{const}$ sitting inside a
raw distribution nudged toward Gaussian --- not the single coherent surface of
\eqref{eq:joint}. This is precisely why exp-34e yields good, calibrated-looking
features but not the log-likelihood identities of \S6: the softmax discards the
constant \emph{and} the interaction lives in a different space than the marginal
constraint.

\subsection{The one-line fix: NPLM on the raw space}
Replace NT-Xent-on-the-sphere with the NPLM ratio on the \emph{same raw} $z$ that
SIGReg controls, using the critic \eqref{eq:critic} (or its distance form). Now
interaction and marginal share one Euclidean geometry, the gauge argument of \S6.3
applies verbatim, and \eqref{eq:marg}--\eqref{eq:joint} hold. The upgrade is a
single structural swap,
\begin{equation}
\text{normalize}+\text{softmax}\ \longrightarrow\ \text{raw}+\text{NPLM},
\end{equation}
turning ``contrastive features that happen to be Gaussian'' into ``an embedding
that is log-likelihood space.''

\subsection{One class, many configurations}
All of the above are instances of a single objective
\begin{equation}
\mathcal{L}
= \mathcal{L}_{\text{int}}\big(g;\ \mathcal{P}\big)
+ \lambda\,\mathcal{L}_{\text{marg}}(z),
\end{equation}
with four orthogonal knobs:
\begin{itemize}
\item \textbf{positives} $\mathcal{P}$: instance (augmented views) $\mid$
supervised (same label) $\mid$ residual ($z-\bar z_y$);
\item \textbf{critic} $g$: cosine $\langle\hat z,\hat z'\rangle/\tau$ $\mid$
bilinear $\langle z,z'\rangle$ $\mid$ neg-distance
$-\tfrac12\|z-z'\|^2$ ($+\,b(x)+b(x')$);
\item \textbf{interaction estimator}: softmax NT-Xent (self-normalized) $\mid$
NPLM $e^{g}-1$ (un-normalized MLE);
\item \textbf{marginal} $\mathcal{L}_{\text{marg}}$: none $\mid$
SIGReg $\to\mathcal{N}(0,I)$ $\mid$ class-wise SIGReg
$\to\mathcal{N}(m_c,\sigma_c^2 I)$.
\end{itemize}
The known recipes are corners of this cube:
\begin{center}
\begin{tabular}{@{}lcccc@{}}
\toprule
recipe & positives & critic & estimator & marginal\\
\midrule
SimCLR & instance & cosine & softmax & none\\
SupCon & supervised & cosine & softmax & none\\
exp-34e hybrid & instance & cosine & softmax & SIGReg\\
34g supcon$+$sigreg & supervised & cosine & softmax & SIGReg\\
LeJEPA-like & pred/target & neg-dist (MSE) & --- & SIGReg\\
\textbf{NPLM geometry} & instance/sup & bilinear/neg-dist & \textbf{NPLM} & SIGReg\\
class-wise NPLM & supervised & neg-dist & NPLM & class-wise SIGReg\\
\bottomrule
\end{tabular}
\end{center}
Making \texttt{train\_simclr\_sigreg} a configurable class along these axes means
every existing exp-34 arm is a special case and the log-likelihood-space arm is one
config flip away --- directly comparable under the same probe protocol.

\subsection{Practical notes for robustness across configurations}
\begin{itemize}
\item \textbf{Stabilize $e^{g}$ by clamping, not by subtracting a max.} It is
tempting to copy the \texttt{sim - sim.max()} guard from \texttt{supcon\_loss},
but that trick is only valid inside a log-\emph{ratio}: in softmax the subtracted
constant cancels, whereas in the NPLM loss there is no enclosing ratio, so
subtracting a constant $c$ shifts the minimizer to $g^\star+c$ and \emph{destroys
the absolute calibration} that is the entire point. Instead clamp the exponent,
$e^{\min(g,\,g_{\max})}$; this tames the blow-up of \S5.1 while leaving $g^\star$
unchanged in the normal range. (Empirically, with clamping the recovered critic
matches the analytic $\operatorname{PMI}$ with correlation $0.99$ and a mean
offset of $\approx 0$; with max-subtraction the offset runs to hundreds.)
\item \textbf{Mask same-class pairs from the reference.} With supervised positives,
off-diagonal same-class pairs are \emph{positives}, not independent draws; they
must be excluded from the $\mathbb{E}_{\mathrm{ref}}[e^{g}]$ sum or they bias the
product-of-marginals estimate --- the NPLM analogue of SupCon's positive masking.
\item \textbf{Drop the bias under class-wise marginals.} When
$\mathcal{L}_{\text{marg}}$ is class-wise SIGReg, the class mean $m_c$ already
supplies the per-point offset, so $b(x)$ can be fixed to $-\tfrac12\|z-m_{y}\|^2$
rather than learned.
\item \textbf{Fold temperature into $g$.} A single scale $\tau$ (or a learnable
scalar on $g$) subsumes the cosine temperature and the NPLM output scale, keeping
the knobs independent and the configurations comparable.
\end{itemize}

\end{document}
